
\documentclass{article}
\usepackage{colortbl}
\usepackage{makecell}
\usepackage{multirow}
\usepackage{supertabular}

\begin{document}

\newcounter{utterance}

\centering \large Interaction Transcript for game `codenames', experiment `association\_easy', episode 0 with Qwen3.6{-}35B{-}A3B{-}FP8{-}without{-}reasoning.
\vspace{24pt}

{ \footnotesize  \setcounter{utterance}{1}
\setlength{\tabcolsep}{0pt}
\begin{supertabular}{c@{$\;$}|p{.15\linewidth}@{}p{.15\linewidth}p{.15\linewidth}p{.15\linewidth}p{.15\linewidth}p{.15\linewidth}}
    \# & $\;$A & \multicolumn{4}{c}{Game Master} & $\;\:$B\\
    \hline

    \theutterance \stepcounter{utterance}  
    & & & \multicolumn{2}{p{0.3\linewidth}}{
        \cellcolor[rgb]{0.9,0.9,0.9}{
            \makecell[{{p{\linewidth}}}]{
                \texttt{\tiny{[GM$|$GM]}}
                {'hidden': {'team': ['czerwony', 'niebieski', 'młotek', 'ksiądz', 'szydło', 'opalony', 'wiertło', 'imam', 'kaznodzieja'], 'innocent': ['miedziogłów', 'czarny', 'siatkówka', 'czołg', 'szczur', 'zapasy', 'kij'], 'opponent': ['trzepaczka', 'płaszcz', 'szczypce', 'talerz', 'sweter', 'koszula', 'durszlak', 't-shirt'], 'assassin': ['wzgórze']}, 'revealed': {'team': {'team': [], 'innocent': [], 'opponent': [], 'assassin': []}, 'opponent': {'team': [], 'innocent': [], 'opponent': [], 'assassin': []}}}
            }
        }
    }
    & & \\ \\

    \theutterance \stepcounter{utterance}  
    & & \multicolumn{4}{p{0.6\linewidth}}{
        \cellcolor[rgb]{0.9,0.9,0.9}{
            \makecell[{{p{\linewidth}}}]{
                \texttt{\tiny{[P1$\langle$GM]}}
                \texttt{Podaj jedną wskazówkę składającą się z pojedynczego słowa, która odnosi się do niektórych słów twojej drużyny z poniższej listy. Możesz wybrać, do ilu i do jakich słów odnosi się twoja wskazówka.} \\
\texttt{Celem jest, aby twój partner odgadł wszystkie słowa twojej drużyny na podstawie podanych wskazówek, zanim drużyna przeciwna odkryje wszystkie swoje. Jeśli twój partner odgadnie słowo 'przeciwnika' lub 'rozpraszające', jego tura zgadywania kończy się i drużyna przeciwna wykonuje turę. Jeśli twój partner odgadnie słowo 'zabójcze', gra kończy się i twoja drużyna natychmiast przegrywa.} \\
\texttt{Im więcej słów odnosi się do twojej wskazówki, tym szybciej możesz wygrać grę, ale tym trudniejsze mogą być skojarzenia do odgadnięcia dla twojego partnera, więc wybieraj słowo wskazówki mądrze. Słowo wskazówki musi być semantycznie powiązane ze słowami docelowymi, nie może być jednym ze słów na listach ani zawierać ich części.} \\
\\ 
\texttt{Zawsze podawaj swoją wskazówkę składającą się z pojedynczego słowa i oddzieloną przecinkami listę powiązanych słów docelowych w następującym formacie i staraj się, aby twoje odpowiedzi były jak najkrótsze, nigdy nie dołączaj żadnego innego tekstu niż jest to wymagane w tym formacie:} \\
\\ 
\texttt{WSKAZÓWKA: <SŁOWO>} \\
\texttt{CELE: <CELE>} \\
\\ 
\texttt{Słowa twojej drużyny to: czerwony, niebieski, młotek, ksiądz, szydło, opalony, wiertło, imam, kaznodzieja.} \\
\texttt{Słowa drużyny przeciwnika to: trzepaczka, płaszcz, szczypce, talerz, sweter, koszula, durszlak, t{-}shirt.} \\
\texttt{Słowa rozpraszające to: miedziogłów, czarny, siatkówka, czołg, szczur, zapasy, kij.} \\
\texttt{Słowa zabójcze to: wzgórze.} \\
            }
        }
    }
    & \\ \\

    \theutterance \stepcounter{utterance}  
    & \multicolumn{4}{p{0.6\linewidth}}{
        \cellcolor[rgb]{0.9,0.9,0.9}{
            \makecell[{{p{\linewidth}}}]{
                \texttt{\tiny{[P1$\rangle$GM]}}
                \texttt{WSKAZÓWKA: narzędzie} \\
\texttt{CELE: młotek, szydło, wiertło} \\
            }
        }
    }
    & & \\ \\

    \theutterance \stepcounter{utterance}  
    & & & \multicolumn{2}{p{0.3\linewidth}}{
        \cellcolor[rgb]{0.9,0.9,0.9}{
            \makecell[{{p{\linewidth}}}]{
                \texttt{\tiny{[GM$|$GM]}}
                \texttt{narzędzie} \\
            }
        }
    }
    & & \\ \\

    \theutterance \stepcounter{utterance}  
    & & & \multicolumn{2}{p{0.3\linewidth}}{
        \cellcolor[rgb]{0.9,0.9,0.9}{
            \makecell[{{p{\linewidth}}}]{
                \texttt{\tiny{[GM$|$GM]}}
                ['młotek', 'szydło', 'wiertło']
            }
        }
    }
    & & \\ \\

    \theutterance \stepcounter{utterance}  
    & & & \multicolumn{2}{p{0.3\linewidth}}{
        \cellcolor[rgb]{0.9,0.9,0.9}{
            \makecell[{{p{\linewidth}}}]{
                \texttt{\tiny{[GM$|$GM]}}
                {'word': 'młotek', 'assignment': 'team'}
            }
        }
    }
    & & \\ \\

    \theutterance \stepcounter{utterance}  
    & & & \multicolumn{2}{p{0.3\linewidth}}{
        \cellcolor[rgb]{0.9,0.9,0.9}{
            \makecell[{{p{\linewidth}}}]{
                \texttt{\tiny{[GM$|$GM]}}
                {'word': 'szydło', 'assignment': 'team'}
            }
        }
    }
    & & \\ \\

    \theutterance \stepcounter{utterance}  
    & & & \multicolumn{2}{p{0.3\linewidth}}{
        \cellcolor[rgb]{0.9,0.9,0.9}{
            \makecell[{{p{\linewidth}}}]{
                \texttt{\tiny{[GM$|$GM]}}
                {'word': 'wiertło', 'assignment': 'team'}
            }
        }
    }
    & & \\ \\

    \theutterance \stepcounter{utterance}  
    & & \multicolumn{4}{p{0.6\linewidth}}{
        \cellcolor[rgb]{0.9,0.9,0.9}{
            \makecell[{{p{\linewidth}}}]{
                \texttt{\tiny{[GM$\rangle$P2]}}
                \texttt{Podaj oddzieloną przecinkami listę maksymalnie 3 słów z poniższej listy, które najlepiej odnoszą się lub są najbardziej związane ze słowem 'narzędzie'. Zawsze rozpoczynaj swoją listę odgadnięć od ZGADNIĘCIE:  i nie dołączaj żadnego innego tekstu w swojej odpowiedzi.} \\
\\ 
\texttt{t{-}shirt, talerz, sweter, imam, szczur, opalony, wiertło, czołg, koszula, siatkówka, szczypce, wzgórze, kaznodzieja, czarny, miedziogłów, czerwony, durszlak, płaszcz, zapasy, niebieski, młotek, kij, ksiądz, trzepaczka, szydło} \\
            }
        }
    }
    & \\ \\

    \theutterance \stepcounter{utterance}  
    & & & \multicolumn{4}{p{0.6\linewidth}}{
        \cellcolor[rgb]{0.9,0.9,0.9}{
            \makecell[{{p{\linewidth}}}]{
                \texttt{\tiny{[GM$\langle$P2]}}
                \texttt{ZGADNIĘCIE: wiertło, szczypce, młotek} \\
            }
        }
    }
     \\ \\

    \theutterance \stepcounter{utterance}  
    & & & \multicolumn{2}{p{0.3\linewidth}}{
        \cellcolor[rgb]{0.9,0.9,0.9}{
            \makecell[{{p{\linewidth}}}]{
                \texttt{\tiny{[GM$|$GM]}}
                ['wiertło', 'szczypce', 'młotek']
            }
        }
    }
    & & \\ \\

    \theutterance \stepcounter{utterance}  
    & & & \multicolumn{2}{p{0.3\linewidth}}{
        \cellcolor[rgb]{0.9,0.9,0.9}{
            \makecell[{{p{\linewidth}}}]{
                \texttt{\tiny{[GM$|$GM]}}
                {'word': 'wiertło', 'assignment': 'team'}
            }
        }
    }
    & & \\ \\

    \theutterance \stepcounter{utterance}  
    & & & \multicolumn{2}{p{0.3\linewidth}}{
        \cellcolor[rgb]{0.9,0.9,0.9}{
            \makecell[{{p{\linewidth}}}]{
                \texttt{\tiny{[GM$|$GM]}}
                {'word': 'wiertło', 'assignment': 'team'}
            }
        }
    }
    & & \\ \\

    \theutterance \stepcounter{utterance}  
    & & & \multicolumn{2}{p{0.3\linewidth}}{
        \cellcolor[rgb]{0.9,0.9,0.9}{
            \makecell[{{p{\linewidth}}}]{
                \texttt{\tiny{[GM$|$GM]}}
                {'word': 'szczypce', 'assignment': 'opponent'}
            }
        }
    }
    & & \\ \\

    \theutterance \stepcounter{utterance}  
    & & & \multicolumn{2}{p{0.3\linewidth}}{
        \cellcolor[rgb]{0.9,0.9,0.9}{
            \makecell[{{p{\linewidth}}}]{
                \texttt{\tiny{[GM$|$GM]}}
                \texttt{szczypce} \\
            }
        }
    }
    & & \\ \\

    \theutterance \stepcounter{utterance}  
    & & & \multicolumn{2}{p{0.3\linewidth}}{
        \cellcolor[rgb]{0.9,0.9,0.9}{
            \makecell[{{p{\linewidth}}}]{
                \texttt{\tiny{[GM$|$GM]}}
                {'word': 'trzepaczka', 'assignment': 'opponent'}
            }
        }
    }
    & & \\ \\

    \theutterance \stepcounter{utterance}  
    & & & \multicolumn{2}{p{0.3\linewidth}}{
        \cellcolor[rgb]{0.9,0.9,0.9}{
            \makecell[{{p{\linewidth}}}]{
                \texttt{\tiny{[GM$|$GM]}}
                {'hidden': {'team': ['czerwony', 'niebieski', 'młotek', 'ksiądz', 'szydło', 'opalony', 'imam', 'kaznodzieja'], 'innocent': ['miedziogłów', 'czarny', 'siatkówka', 'czołg', 'szczur', 'zapasy', 'kij'], 'opponent': ['płaszcz', 'talerz', 'sweter', 'koszula', 'durszlak', 't-shirt'], 'assassin': ['wzgórze']}, 'revealed': {'team': {'team': ['wiertło'], 'innocent': [], 'opponent': ['szczypce'], 'assassin': []}, 'opponent': {'team': [], 'innocent': [], 'opponent': ['trzepaczka'], 'assassin': []}}}
            }
        }
    }
    & & \\ \\

    \theutterance \stepcounter{utterance}  
    & & \multicolumn{4}{p{0.6\linewidth}}{
        \cellcolor[rgb]{0.9,0.9,0.9}{
            \makecell[{{p{\linewidth}}}]{
                \texttt{\tiny{[P1$\langle$GM]}}
                \texttt{Słowo wiertło zostało odgadnięte prawidłowo. Słowo szczypce zostało odgadnięte, ale to słowo przeciwnik. Tura twojego kolegi z drużyny się tam skończyła.} \\
\texttt{Teraz podaj kolejną wskazówkę odnoszącą się do niektórych z pozostałych słów twojego zespołu oraz listę powiązanych słów docelowych. Pamiętaj, aby rozpocząć wskazówkę od 'WSKAZÓWKA: ', wstawić nową linię i rozpocząć listę słów docelowych oddzielonych przecinkami od 'CELE: '. Uwaga: niektóre słowa zostały usunięte z list w porównaniu do poprzednich żądań.} \\
\\ 
\texttt{Twoje pozostałe słowa zespołu to: czerwony, niebieski, młotek, ksiądz, szydło, opalony, imam, kaznodzieja.} \\
\texttt{Pozostałe słowa twojego przeciwnika to: trzepaczka, płaszcz, talerz, sweter, koszula, durszlak, t{-}shirt.} \\
\texttt{Pozostałe słowa rozpraszające to: miedziogłów, czarny, siatkówka, czołg, szczur, zapasy, kij.} \\
\texttt{Pozostałe słowa zabójcy to: wzgórze.} \\
            }
        }
    }
    & \\ \\

    \theutterance \stepcounter{utterance}  
    & \multicolumn{4}{p{0.6\linewidth}}{
        \cellcolor[rgb]{0.9,0.9,0.9}{
            \makecell[{{p{\linewidth}}}]{
                \texttt{\tiny{[P1$\rangle$GM]}}
                \texttt{WSKAZÓWKA: duchowny} \\
\texttt{CELE: ksiądz, imam, kaznodzieja} \\
            }
        }
    }
    & & \\ \\

    \theutterance \stepcounter{utterance}  
    & & & \multicolumn{2}{p{0.3\linewidth}}{
        \cellcolor[rgb]{0.9,0.9,0.9}{
            \makecell[{{p{\linewidth}}}]{
                \texttt{\tiny{[GM$|$GM]}}
                \texttt{duchowny} \\
            }
        }
    }
    & & \\ \\

    \theutterance \stepcounter{utterance}  
    & & & \multicolumn{2}{p{0.3\linewidth}}{
        \cellcolor[rgb]{0.9,0.9,0.9}{
            \makecell[{{p{\linewidth}}}]{
                \texttt{\tiny{[GM$|$GM]}}
                ['ksiądz', 'imam', 'kaznodzieja']
            }
        }
    }
    & & \\ \\

    \theutterance \stepcounter{utterance}  
    & & & \multicolumn{2}{p{0.3\linewidth}}{
        \cellcolor[rgb]{0.9,0.9,0.9}{
            \makecell[{{p{\linewidth}}}]{
                \texttt{\tiny{[GM$|$GM]}}
                {'word': 'ksiądz', 'assignment': 'team'}
            }
        }
    }
    & & \\ \\

    \theutterance \stepcounter{utterance}  
    & & & \multicolumn{2}{p{0.3\linewidth}}{
        \cellcolor[rgb]{0.9,0.9,0.9}{
            \makecell[{{p{\linewidth}}}]{
                \texttt{\tiny{[GM$|$GM]}}
                {'word': 'imam', 'assignment': 'team'}
            }
        }
    }
    & & \\ \\

    \theutterance \stepcounter{utterance}  
    & & & \multicolumn{2}{p{0.3\linewidth}}{
        \cellcolor[rgb]{0.9,0.9,0.9}{
            \makecell[{{p{\linewidth}}}]{
                \texttt{\tiny{[GM$|$GM]}}
                {'word': 'kaznodzieja', 'assignment': 'team'}
            }
        }
    }
    & & \\ \\

    \theutterance \stepcounter{utterance}  
    & & \multicolumn{4}{p{0.6\linewidth}}{
        \cellcolor[rgb]{0.9,0.9,0.9}{
            \makecell[{{p{\linewidth}}}]{
                \texttt{\tiny{[GM$\rangle$P2]}}
                \texttt{Słowo wiertło zostało odgadnięte prawidłowo. Słowo szczypce zostało odgadnięte, ale to słowo przeciwnik. Twoja tura się tam skończyła.} \\
\texttt{Teraz podaj kolejną listę oddzieloną przecinkami, zawierającą co najmniej 1 i maksymalnie 3 słów z poniższej listy słów, które najlepiej odnoszą się lub są najbardziej związane ze słowem 'duchowny'. Pamiętaj, aby rozpocząć odpowiedź od ZGADNIĘCIE: . Uwaga: niektóre słowa zostały usunięte z listy w porównaniu do poprzednich próśb.} \\
\\ 
\texttt{t{-}shirt, talerz, sweter, imam, szczur, opalony, czołg, koszula, siatkówka, wzgórze, kaznodzieja, czarny, miedziogłów, czerwony, durszlak, płaszcz, zapasy, niebieski, młotek, kij, ksiądz, szydło} \\
            }
        }
    }
    & \\ \\

    \theutterance \stepcounter{utterance}  
    & & & \multicolumn{4}{p{0.6\linewidth}}{
        \cellcolor[rgb]{0.9,0.9,0.9}{
            \makecell[{{p{\linewidth}}}]{
                \texttt{\tiny{[GM$\langle$P2]}}
                \texttt{ZGADNIĘCIE: imam, kaznodzieja, ksiądz} \\
            }
        }
    }
     \\ \\

    \theutterance \stepcounter{utterance}  
    & & & \multicolumn{2}{p{0.3\linewidth}}{
        \cellcolor[rgb]{0.9,0.9,0.9}{
            \makecell[{{p{\linewidth}}}]{
                \texttt{\tiny{[GM$|$GM]}}
                ['imam', 'kaznodzieja', 'ksiądz']
            }
        }
    }
    & & \\ \\

    \theutterance \stepcounter{utterance}  
    & & & \multicolumn{2}{p{0.3\linewidth}}{
        \cellcolor[rgb]{0.9,0.9,0.9}{
            \makecell[{{p{\linewidth}}}]{
                \texttt{\tiny{[GM$|$GM]}}
                {'word': 'imam', 'assignment': 'team'}
            }
        }
    }
    & & \\ \\

    \theutterance \stepcounter{utterance}  
    & & & \multicolumn{2}{p{0.3\linewidth}}{
        \cellcolor[rgb]{0.9,0.9,0.9}{
            \makecell[{{p{\linewidth}}}]{
                \texttt{\tiny{[GM$|$GM]}}
                {'word': 'imam', 'assignment': 'team'}
            }
        }
    }
    & & \\ \\

    \theutterance \stepcounter{utterance}  
    & & & \multicolumn{2}{p{0.3\linewidth}}{
        \cellcolor[rgb]{0.9,0.9,0.9}{
            \makecell[{{p{\linewidth}}}]{
                \texttt{\tiny{[GM$|$GM]}}
                {'word': 'kaznodzieja', 'assignment': 'team'}
            }
        }
    }
    & & \\ \\

    \theutterance \stepcounter{utterance}  
    & & & \multicolumn{2}{p{0.3\linewidth}}{
        \cellcolor[rgb]{0.9,0.9,0.9}{
            \makecell[{{p{\linewidth}}}]{
                \texttt{\tiny{[GM$|$GM]}}
                {'word': 'kaznodzieja', 'assignment': 'team'}
            }
        }
    }
    & & \\ \\

    \theutterance \stepcounter{utterance}  
    & & & \multicolumn{2}{p{0.3\linewidth}}{
        \cellcolor[rgb]{0.9,0.9,0.9}{
            \makecell[{{p{\linewidth}}}]{
                \texttt{\tiny{[GM$|$GM]}}
                {'word': 'ksiądz', 'assignment': 'team'}
            }
        }
    }
    & & \\ \\

    \theutterance \stepcounter{utterance}  
    & & & \multicolumn{2}{p{0.3\linewidth}}{
        \cellcolor[rgb]{0.9,0.9,0.9}{
            \makecell[{{p{\linewidth}}}]{
                \texttt{\tiny{[GM$|$GM]}}
                {'word': 'ksiądz', 'assignment': 'team'}
            }
        }
    }
    & & \\ \\

    \theutterance \stepcounter{utterance}  
    & & & \multicolumn{2}{p{0.3\linewidth}}{
        \cellcolor[rgb]{0.9,0.9,0.9}{
            \makecell[{{p{\linewidth}}}]{
                \texttt{\tiny{[GM$|$GM]}}
                {'word': 'sweter', 'assignment': 'opponent'}
            }
        }
    }
    & & \\ \\

    \theutterance \stepcounter{utterance}  
    & & & \multicolumn{2}{p{0.3\linewidth}}{
        \cellcolor[rgb]{0.9,0.9,0.9}{
            \makecell[{{p{\linewidth}}}]{
                \texttt{\tiny{[GM$|$GM]}}
                {'hidden': {'team': ['czerwony', 'niebieski', 'młotek', 'szydło', 'opalony'], 'innocent': ['miedziogłów', 'czarny', 'siatkówka', 'czołg', 'szczur', 'zapasy', 'kij'], 'opponent': ['płaszcz', 'talerz', 'koszula', 'durszlak', 't-shirt'], 'assassin': ['wzgórze']}, 'revealed': {'team': {'team': ['wiertło', 'imam', 'kaznodzieja', 'ksiądz'], 'innocent': [], 'opponent': ['szczypce'], 'assassin': []}, 'opponent': {'team': [], 'innocent': [], 'opponent': ['trzepaczka', 'sweter'], 'assassin': []}}}
            }
        }
    }
    & & \\ \\

    \theutterance \stepcounter{utterance}  
    & & \multicolumn{4}{p{0.6\linewidth}}{
        \cellcolor[rgb]{0.9,0.9,0.9}{
            \makecell[{{p{\linewidth}}}]{
                \texttt{\tiny{[P1$\langle$GM]}}
                \texttt{Słowa imam, kaznodzieja, ksiądz zostały odgadnięte prawidłowo. Tura twojego kolegi z drużyny się tam skończyła.} \\
\texttt{Teraz podaj kolejną wskazówkę odnoszącą się do niektórych z pozostałych słów twojego zespołu oraz listę powiązanych słów docelowych. Pamiętaj, aby rozpocząć wskazówkę od 'WSKAZÓWKA: ', wstawić nową linię i rozpocząć listę słów docelowych oddzielonych przecinkami od 'CELE: '. Uwaga: niektóre słowa zostały usunięte z list w porównaniu do poprzednich żądań.} \\
\\ 
\texttt{Twoje pozostałe słowa zespołu to: czerwony, niebieski, młotek, szydło, opalony.} \\
\texttt{Pozostałe słowa twojego przeciwnika to: płaszcz, talerz, sweter, koszula, durszlak, t{-}shirt.} \\
\texttt{Pozostałe słowa rozpraszające to: miedziogłów, czarny, siatkówka, czołg, szczur, zapasy, kij.} \\
\texttt{Pozostałe słowa zabójcy to: wzgórze.} \\
            }
        }
    }
    & \\ \\

    \theutterance \stepcounter{utterance}  
    & \multicolumn{4}{p{0.6\linewidth}}{
        \cellcolor[rgb]{0.9,0.9,0.9}{
            \makecell[{{p{\linewidth}}}]{
                \texttt{\tiny{[P1$\rangle$GM]}}
                \texttt{WSKAZÓWKA: kolor} \\
\texttt{CELE: czerwony, niebieski} \\
            }
        }
    }
    & & \\ \\

    \theutterance \stepcounter{utterance}  
    & & & \multicolumn{2}{p{0.3\linewidth}}{
        \cellcolor[rgb]{0.9,0.9,0.9}{
            \makecell[{{p{\linewidth}}}]{
                \texttt{\tiny{[GM$|$GM]}}
                \texttt{kolor} \\
            }
        }
    }
    & & \\ \\

    \theutterance \stepcounter{utterance}  
    & & & \multicolumn{2}{p{0.3\linewidth}}{
        \cellcolor[rgb]{0.9,0.9,0.9}{
            \makecell[{{p{\linewidth}}}]{
                \texttt{\tiny{[GM$|$GM]}}
                ['czerwony', 'niebieski']
            }
        }
    }
    & & \\ \\

    \theutterance \stepcounter{utterance}  
    & & & \multicolumn{2}{p{0.3\linewidth}}{
        \cellcolor[rgb]{0.9,0.9,0.9}{
            \makecell[{{p{\linewidth}}}]{
                \texttt{\tiny{[GM$|$GM]}}
                {'word': 'czerwony', 'assignment': 'team'}
            }
        }
    }
    & & \\ \\

    \theutterance \stepcounter{utterance}  
    & & & \multicolumn{2}{p{0.3\linewidth}}{
        \cellcolor[rgb]{0.9,0.9,0.9}{
            \makecell[{{p{\linewidth}}}]{
                \texttt{\tiny{[GM$|$GM]}}
                {'word': 'niebieski', 'assignment': 'team'}
            }
        }
    }
    & & \\ \\

    \theutterance \stepcounter{utterance}  
    & & \multicolumn{4}{p{0.6\linewidth}}{
        \cellcolor[rgb]{0.9,0.9,0.9}{
            \makecell[{{p{\linewidth}}}]{
                \texttt{\tiny{[GM$\rangle$P2]}}
                \texttt{Słowa imam, kaznodzieja, ksiądz zostały odgadnięte prawidłowo. Twoja tura się tam skończyła.} \\
\texttt{Teraz podaj kolejną listę oddzieloną przecinkami, zawierającą co najmniej 1 i maksymalnie 2 słów z poniższej listy słów, które najlepiej odnoszą się lub są najbardziej związane ze słowem 'kolor'. Pamiętaj, aby rozpocząć odpowiedź od ZGADNIĘCIE: . Uwaga: niektóre słowa zostały usunięte z listy w porównaniu do poprzednich próśb.} \\
\\ 
\texttt{t{-}shirt, talerz, szczur, opalony, czołg, koszula, siatkówka, wzgórze, czarny, miedziogłów, czerwony, durszlak, płaszcz, zapasy, niebieski, młotek, kij, szydło} \\
            }
        }
    }
    & \\ \\

    \theutterance \stepcounter{utterance}  
    & & & \multicolumn{4}{p{0.6\linewidth}}{
        \cellcolor[rgb]{0.9,0.9,0.9}{
            \makecell[{{p{\linewidth}}}]{
                \texttt{\tiny{[GM$\langle$P2]}}
                \texttt{ZGADNIĘCIE: czarny, czerwony} \\
            }
        }
    }
     \\ \\

    \theutterance \stepcounter{utterance}  
    & & & \multicolumn{2}{p{0.3\linewidth}}{
        \cellcolor[rgb]{0.9,0.9,0.9}{
            \makecell[{{p{\linewidth}}}]{
                \texttt{\tiny{[GM$|$GM]}}
                ['czarny', 'czerwony']
            }
        }
    }
    & & \\ \\

    \theutterance \stepcounter{utterance}  
    & & & \multicolumn{2}{p{0.3\linewidth}}{
        \cellcolor[rgb]{0.9,0.9,0.9}{
            \makecell[{{p{\linewidth}}}]{
                \texttt{\tiny{[GM$|$GM]}}
                {'word': 'czarny', 'assignment': 'innocent'}
            }
        }
    }
    & & \\ \\

    \theutterance \stepcounter{utterance}  
    & & & \multicolumn{2}{p{0.3\linewidth}}{
        \cellcolor[rgb]{0.9,0.9,0.9}{
            \makecell[{{p{\linewidth}}}]{
                \texttt{\tiny{[GM$|$GM]}}
                \texttt{czarny} \\
            }
        }
    }
    & & \\ \\

    \theutterance \stepcounter{utterance}  
    & & & \multicolumn{2}{p{0.3\linewidth}}{
        \cellcolor[rgb]{0.9,0.9,0.9}{
            \makecell[{{p{\linewidth}}}]{
                \texttt{\tiny{[GM$|$GM]}}
                {'word': 'koszula', 'assignment': 'opponent'}
            }
        }
    }
    & & \\ \\

    \theutterance \stepcounter{utterance}  
    & & & \multicolumn{2}{p{0.3\linewidth}}{
        \cellcolor[rgb]{0.9,0.9,0.9}{
            \makecell[{{p{\linewidth}}}]{
                \texttt{\tiny{[GM$|$GM]}}
                {'hidden': {'team': ['czerwony', 'niebieski', 'młotek', 'szydło', 'opalony'], 'innocent': ['miedziogłów', 'siatkówka', 'czołg', 'szczur', 'zapasy', 'kij'], 'opponent': ['płaszcz', 'talerz', 'durszlak', 't-shirt'], 'assassin': ['wzgórze']}, 'revealed': {'team': {'team': ['wiertło', 'imam', 'kaznodzieja', 'ksiądz'], 'innocent': ['czarny'], 'opponent': ['szczypce'], 'assassin': []}, 'opponent': {'team': [], 'innocent': [], 'opponent': ['trzepaczka', 'sweter', 'koszula'], 'assassin': []}}}
            }
        }
    }
    & & \\ \\

    \theutterance \stepcounter{utterance}  
    & & \multicolumn{4}{p{0.6\linewidth}}{
        \cellcolor[rgb]{0.9,0.9,0.9}{
            \makecell[{{p{\linewidth}}}]{
                \texttt{\tiny{[P1$\langle$GM]}}
                \texttt{Słowo czarny zostało odgadnięte, ale to słowo niewinny. Tura twojego kolegi z drużyny się tam skończyła.} \\
\texttt{Teraz podaj kolejną wskazówkę odnoszącą się do niektórych z pozostałych słów twojego zespołu oraz listę powiązanych słów docelowych. Pamiętaj, aby rozpocząć wskazówkę od 'WSKAZÓWKA: ', wstawić nową linię i rozpocząć listę słów docelowych oddzielonych przecinkami od 'CELE: '. Uwaga: niektóre słowa zostały usunięte z list w porównaniu do poprzednich żądań.} \\
\\ 
\texttt{Twoje pozostałe słowa zespołu to: czerwony, niebieski, młotek, szydło, opalony.} \\
\texttt{Pozostałe słowa twojego przeciwnika to: płaszcz, talerz, koszula, durszlak, t{-}shirt.} \\
\texttt{Pozostałe słowa rozpraszające to: miedziogłów, siatkówka, czołg, szczur, zapasy, kij.} \\
\texttt{Pozostałe słowa zabójcy to: wzgórze.} \\
            }
        }
    }
    & \\ \\

    \theutterance \stepcounter{utterance}  
    & \multicolumn{4}{p{0.6\linewidth}}{
        \cellcolor[rgb]{0.9,0.9,0.9}{
            \makecell[{{p{\linewidth}}}]{
                \texttt{\tiny{[P1$\rangle$GM]}}
                \texttt{WSKAZÓWKA: ubranie} \\
\texttt{CELE: czerwony, niebieski} \\
            }
        }
    }
    & & \\ \\

    \theutterance \stepcounter{utterance}  
    & & & \multicolumn{2}{p{0.3\linewidth}}{
        \cellcolor[rgb]{0.9,0.9,0.9}{
            \makecell[{{p{\linewidth}}}]{
                \texttt{\tiny{[GM$|$GM]}}
                \texttt{ubranie} \\
            }
        }
    }
    & & \\ \\

    \theutterance \stepcounter{utterance}  
    & & & \multicolumn{2}{p{0.3\linewidth}}{
        \cellcolor[rgb]{0.9,0.9,0.9}{
            \makecell[{{p{\linewidth}}}]{
                \texttt{\tiny{[GM$|$GM]}}
                ['czerwony', 'niebieski']
            }
        }
    }
    & & \\ \\

    \theutterance \stepcounter{utterance}  
    & & & \multicolumn{2}{p{0.3\linewidth}}{
        \cellcolor[rgb]{0.9,0.9,0.9}{
            \makecell[{{p{\linewidth}}}]{
                \texttt{\tiny{[GM$|$GM]}}
                {'word': 'czerwony', 'assignment': 'team'}
            }
        }
    }
    & & \\ \\

    \theutterance \stepcounter{utterance}  
    & & & \multicolumn{2}{p{0.3\linewidth}}{
        \cellcolor[rgb]{0.9,0.9,0.9}{
            \makecell[{{p{\linewidth}}}]{
                \texttt{\tiny{[GM$|$GM]}}
                {'word': 'niebieski', 'assignment': 'team'}
            }
        }
    }
    & & \\ \\

    \theutterance \stepcounter{utterance}  
    & & \multicolumn{4}{p{0.6\linewidth}}{
        \cellcolor[rgb]{0.9,0.9,0.9}{
            \makecell[{{p{\linewidth}}}]{
                \texttt{\tiny{[GM$\rangle$P2]}}
                \texttt{Słowo czarny zostało odgadnięte, ale to słowo niewinny. Twoja tura się tam skończyła.} \\
\texttt{Teraz podaj kolejną listę oddzieloną przecinkami, zawierającą co najmniej 1 i maksymalnie 2 słów z poniższej listy słów, które najlepiej odnoszą się lub są najbardziej związane ze słowem 'ubranie'. Pamiętaj, aby rozpocząć odpowiedź od ZGADNIĘCIE: . Uwaga: niektóre słowa zostały usunięte z listy w porównaniu do poprzednich próśb.} \\
\\ 
\texttt{t{-}shirt, talerz, szczur, opalony, czołg, siatkówka, wzgórze, miedziogłów, czerwony, durszlak, płaszcz, zapasy, niebieski, młotek, kij, szydło} \\
            }
        }
    }
    & \\ \\

    \theutterance \stepcounter{utterance}  
    & & & \multicolumn{4}{p{0.6\linewidth}}{
        \cellcolor[rgb]{0.9,0.9,0.9}{
            \makecell[{{p{\linewidth}}}]{
                \texttt{\tiny{[GM$\langle$P2]}}
                \texttt{ZGADNIĘCIE: t{-}shirt, płaszcz} \\
            }
        }
    }
     \\ \\

    \theutterance \stepcounter{utterance}  
    & & & \multicolumn{2}{p{0.3\linewidth}}{
        \cellcolor[rgb]{0.9,0.9,0.9}{
            \makecell[{{p{\linewidth}}}]{
                \texttt{\tiny{[GM$|$GM]}}
                ['t-shirt', 'płaszcz']
            }
        }
    }
    & & \\ \\

    \theutterance \stepcounter{utterance}  
    & & & \multicolumn{2}{p{0.3\linewidth}}{
        \cellcolor[rgb]{0.9,0.9,0.9}{
            \makecell[{{p{\linewidth}}}]{
                \texttt{\tiny{[GM$|$GM]}}
                {'word': 't-shirt', 'assignment': 'opponent'}
            }
        }
    }
    & & \\ \\

    \theutterance \stepcounter{utterance}  
    & & & \multicolumn{2}{p{0.3\linewidth}}{
        \cellcolor[rgb]{0.9,0.9,0.9}{
            \makecell[{{p{\linewidth}}}]{
                \texttt{\tiny{[GM$|$GM]}}
                \texttt{t{-}shirt} \\
            }
        }
    }
    & & \\ \\

    \theutterance \stepcounter{utterance}  
    & & & \multicolumn{2}{p{0.3\linewidth}}{
        \cellcolor[rgb]{0.9,0.9,0.9}{
            \makecell[{{p{\linewidth}}}]{
                \texttt{\tiny{[GM$|$GM]}}
                {'word': 'durszlak', 'assignment': 'opponent'}
            }
        }
    }
    & & \\ \\

    \theutterance \stepcounter{utterance}  
    & & & \multicolumn{2}{p{0.3\linewidth}}{
        \cellcolor[rgb]{0.9,0.9,0.9}{
            \makecell[{{p{\linewidth}}}]{
                \texttt{\tiny{[GM$|$GM]}}
                {'hidden': {'team': ['czerwony', 'niebieski', 'młotek', 'szydło', 'opalony'], 'innocent': ['miedziogłów', 'siatkówka', 'czołg', 'szczur', 'zapasy', 'kij'], 'opponent': ['płaszcz', 'talerz'], 'assassin': ['wzgórze']}, 'revealed': {'team': {'team': ['wiertło', 'imam', 'kaznodzieja', 'ksiądz'], 'innocent': ['czarny'], 'opponent': ['szczypce', 't-shirt'], 'assassin': []}, 'opponent': {'team': [], 'innocent': [], 'opponent': ['trzepaczka', 'sweter', 'koszula', 'durszlak'], 'assassin': []}}}
            }
        }
    }
    & & \\ \\

    \theutterance \stepcounter{utterance}  
    & & \multicolumn{4}{p{0.6\linewidth}}{
        \cellcolor[rgb]{0.9,0.9,0.9}{
            \makecell[{{p{\linewidth}}}]{
                \texttt{\tiny{[P1$\langle$GM]}}
                \texttt{Słowo t{-}shirt zostało odgadnięte, ale to słowo przeciwnik. Tura twojego kolegi z drużyny się tam skończyła.} \\
\texttt{Teraz podaj kolejną wskazówkę odnoszącą się do niektórych z pozostałych słów twojego zespołu oraz listę powiązanych słów docelowych. Pamiętaj, aby rozpocząć wskazówkę od 'WSKAZÓWKA: ', wstawić nową linię i rozpocząć listę słów docelowych oddzielonych przecinkami od 'CELE: '. Uwaga: niektóre słowa zostały usunięte z list w porównaniu do poprzednich żądań.} \\
\\ 
\texttt{Twoje pozostałe słowa zespołu to: czerwony, niebieski, młotek, szydło, opalony.} \\
\texttt{Pozostałe słowa twojego przeciwnika to: płaszcz, talerz, durszlak.} \\
\texttt{Pozostałe słowa rozpraszające to: miedziogłów, siatkówka, czołg, szczur, zapasy, kij.} \\
\texttt{Pozostałe słowa zabójcy to: wzgórze.} \\
            }
        }
    }
    & \\ \\

    \theutterance \stepcounter{utterance}  
    & \multicolumn{4}{p{0.6\linewidth}}{
        \cellcolor[rgb]{0.9,0.9,0.9}{
            \makecell[{{p{\linewidth}}}]{
                \texttt{\tiny{[P1$\rangle$GM]}}
                \texttt{WSKAZÓWKA: narzędzie} \\
\texttt{CELE: młotek, szydło} \\
            }
        }
    }
    & & \\ \\

    \theutterance \stepcounter{utterance}  
    & & & \multicolumn{2}{p{0.3\linewidth}}{
        \cellcolor[rgb]{0.9,0.9,0.9}{
            \makecell[{{p{\linewidth}}}]{
                \texttt{\tiny{[GM$|$GM]}}
                \texttt{narzędzie} \\
            }
        }
    }
    & & \\ \\

    \theutterance \stepcounter{utterance}  
    & & & \multicolumn{2}{p{0.3\linewidth}}{
        \cellcolor[rgb]{0.9,0.9,0.9}{
            \makecell[{{p{\linewidth}}}]{
                \texttt{\tiny{[GM$|$GM]}}
                ['młotek', 'szydło']
            }
        }
    }
    & & \\ \\

    \theutterance \stepcounter{utterance}  
    & & & \multicolumn{2}{p{0.3\linewidth}}{
        \cellcolor[rgb]{0.9,0.9,0.9}{
            \makecell[{{p{\linewidth}}}]{
                \texttt{\tiny{[GM$|$GM]}}
                {'word': 'młotek', 'assignment': 'team'}
            }
        }
    }
    & & \\ \\

    \theutterance \stepcounter{utterance}  
    & & & \multicolumn{2}{p{0.3\linewidth}}{
        \cellcolor[rgb]{0.9,0.9,0.9}{
            \makecell[{{p{\linewidth}}}]{
                \texttt{\tiny{[GM$|$GM]}}
                {'word': 'szydło', 'assignment': 'team'}
            }
        }
    }
    & & \\ \\

    \theutterance \stepcounter{utterance}  
    & & \multicolumn{4}{p{0.6\linewidth}}{
        \cellcolor[rgb]{0.9,0.9,0.9}{
            \makecell[{{p{\linewidth}}}]{
                \texttt{\tiny{[GM$\rangle$P2]}}
                \texttt{Słowo t{-}shirt zostało odgadnięte, ale to słowo przeciwnik. Twoja tura się tam skończyła.} \\
\texttt{Teraz podaj kolejną listę oddzieloną przecinkami, zawierającą co najmniej 1 i maksymalnie 2 słów z poniższej listy słów, które najlepiej odnoszą się lub są najbardziej związane ze słowem 'narzędzie'. Pamiętaj, aby rozpocząć odpowiedź od ZGADNIĘCIE: . Uwaga: niektóre słowa zostały usunięte z listy w porównaniu do poprzednich próśb.} \\
\\ 
\texttt{talerz, szczur, opalony, czołg, siatkówka, wzgórze, miedziogłów, czerwony, płaszcz, zapasy, niebieski, młotek, kij, szydło} \\
            }
        }
    }
    & \\ \\

    \theutterance \stepcounter{utterance}  
    & & & \multicolumn{4}{p{0.6\linewidth}}{
        \cellcolor[rgb]{0.9,0.9,0.9}{
            \makecell[{{p{\linewidth}}}]{
                \texttt{\tiny{[GM$\langle$P2]}}
                \texttt{ZGADNIĘCIE: młotek, kij} \\
            }
        }
    }
     \\ \\

    \theutterance \stepcounter{utterance}  
    & & & \multicolumn{2}{p{0.3\linewidth}}{
        \cellcolor[rgb]{0.9,0.9,0.9}{
            \makecell[{{p{\linewidth}}}]{
                \texttt{\tiny{[GM$|$GM]}}
                ['młotek', 'kij']
            }
        }
    }
    & & \\ \\

    \theutterance \stepcounter{utterance}  
    & & & \multicolumn{2}{p{0.3\linewidth}}{
        \cellcolor[rgb]{0.9,0.9,0.9}{
            \makecell[{{p{\linewidth}}}]{
                \texttt{\tiny{[GM$|$GM]}}
                {'word': 'młotek', 'assignment': 'team'}
            }
        }
    }
    & & \\ \\

    \theutterance \stepcounter{utterance}  
    & & & \multicolumn{2}{p{0.3\linewidth}}{
        \cellcolor[rgb]{0.9,0.9,0.9}{
            \makecell[{{p{\linewidth}}}]{
                \texttt{\tiny{[GM$|$GM]}}
                {'word': 'młotek', 'assignment': 'team'}
            }
        }
    }
    & & \\ \\

    \theutterance \stepcounter{utterance}  
    & & & \multicolumn{2}{p{0.3\linewidth}}{
        \cellcolor[rgb]{0.9,0.9,0.9}{
            \makecell[{{p{\linewidth}}}]{
                \texttt{\tiny{[GM$|$GM]}}
                {'word': 'kij', 'assignment': 'innocent'}
            }
        }
    }
    & & \\ \\

    \theutterance \stepcounter{utterance}  
    & & & \multicolumn{2}{p{0.3\linewidth}}{
        \cellcolor[rgb]{0.9,0.9,0.9}{
            \makecell[{{p{\linewidth}}}]{
                \texttt{\tiny{[GM$|$GM]}}
                \texttt{kij} \\
            }
        }
    }
    & & \\ \\

    \theutterance \stepcounter{utterance}  
    & & & \multicolumn{2}{p{0.3\linewidth}}{
        \cellcolor[rgb]{0.9,0.9,0.9}{
            \makecell[{{p{\linewidth}}}]{
                \texttt{\tiny{[GM$|$GM]}}
                {'word': 'talerz', 'assignment': 'opponent'}
            }
        }
    }
    & & \\ \\

    \theutterance \stepcounter{utterance}  
    & & & \multicolumn{2}{p{0.3\linewidth}}{
        \cellcolor[rgb]{0.9,0.9,0.9}{
            \makecell[{{p{\linewidth}}}]{
                \texttt{\tiny{[GM$|$GM]}}
                {'hidden': {'team': ['czerwony', 'niebieski', 'szydło', 'opalony'], 'innocent': ['miedziogłów', 'siatkówka', 'czołg', 'szczur', 'zapasy'], 'opponent': ['płaszcz'], 'assassin': ['wzgórze']}, 'revealed': {'team': {'team': ['wiertło', 'imam', 'kaznodzieja', 'ksiądz', 'młotek'], 'innocent': ['czarny', 'kij'], 'opponent': ['szczypce', 't-shirt'], 'assassin': []}, 'opponent': {'team': [], 'innocent': [], 'opponent': ['trzepaczka', 'sweter', 'koszula', 'durszlak', 'talerz'], 'assassin': []}}}
            }
        }
    }
    & & \\ \\

    \theutterance \stepcounter{utterance}  
    & & \multicolumn{4}{p{0.6\linewidth}}{
        \cellcolor[rgb]{0.9,0.9,0.9}{
            \makecell[{{p{\linewidth}}}]{
                \texttt{\tiny{[P1$\langle$GM]}}
                \texttt{Słowo młotek zostało odgadnięte prawidłowo. Słowo kij zostało odgadnięte, ale to słowo niewinny. Tura twojego kolegi z drużyny się tam skończyła.} \\
\texttt{Teraz podaj kolejną wskazówkę odnoszącą się do niektórych z pozostałych słów twojego zespołu oraz listę powiązanych słów docelowych. Pamiętaj, aby rozpocząć wskazówkę od 'WSKAZÓWKA: ', wstawić nową linię i rozpocząć listę słów docelowych oddzielonych przecinkami od 'CELE: '. Uwaga: niektóre słowa zostały usunięte z list w porównaniu do poprzednich żądań.} \\
\\ 
\texttt{Twoje pozostałe słowa zespołu to: czerwony, niebieski, szydło, opalony.} \\
\texttt{Pozostałe słowa twojego przeciwnika to: płaszcz, talerz.} \\
\texttt{Pozostałe słowa rozpraszające to: miedziogłów, siatkówka, czołg, szczur, zapasy.} \\
\texttt{Pozostałe słowa zabójcy to: wzgórze.} \\
            }
        }
    }
    & \\ \\

    \theutterance \stepcounter{utterance}  
    & \multicolumn{4}{p{0.6\linewidth}}{
        \cellcolor[rgb]{0.9,0.9,0.9}{
            \makecell[{{p{\linewidth}}}]{
                \texttt{\tiny{[P1$\rangle$GM]}}
                \texttt{WSKAZÓWKA: narzędzie} \\
\texttt{CELE: szydło} \\
            }
        }
    }
    & & \\ \\

    \theutterance \stepcounter{utterance}  
    & & & \multicolumn{2}{p{0.3\linewidth}}{
        \cellcolor[rgb]{0.9,0.9,0.9}{
            \makecell[{{p{\linewidth}}}]{
                \texttt{\tiny{[GM$|$GM]}}
                \texttt{narzędzie} \\
            }
        }
    }
    & & \\ \\

    \theutterance \stepcounter{utterance}  
    & & & \multicolumn{2}{p{0.3\linewidth}}{
        \cellcolor[rgb]{0.9,0.9,0.9}{
            \makecell[{{p{\linewidth}}}]{
                \texttt{\tiny{[GM$|$GM]}}
                ['szydło']
            }
        }
    }
    & & \\ \\

    \theutterance \stepcounter{utterance}  
    & & & \multicolumn{2}{p{0.3\linewidth}}{
        \cellcolor[rgb]{0.9,0.9,0.9}{
            \makecell[{{p{\linewidth}}}]{
                \texttt{\tiny{[GM$|$GM]}}
                {'word': 'szydło', 'assignment': 'team'}
            }
        }
    }
    & & \\ \\

    \theutterance \stepcounter{utterance}  
    & & \multicolumn{4}{p{0.6\linewidth}}{
        \cellcolor[rgb]{0.9,0.9,0.9}{
            \makecell[{{p{\linewidth}}}]{
                \texttt{\tiny{[GM$\rangle$P2]}}
                \texttt{Słowo młotek zostało odgadnięte prawidłowo. Słowo kij zostało odgadnięte, ale to słowo niewinny. Twoja tura się tam skończyła.} \\
\texttt{Teraz podaj kolejną listę oddzieloną przecinkami, zawierającą co najmniej 1 i maksymalnie 1 słów z poniższej listy słów, które najlepiej odnoszą się lub są najbardziej związane ze słowem 'narzędzie'. Pamiętaj, aby rozpocząć odpowiedź od ZGADNIĘCIE: . Uwaga: niektóre słowa zostały usunięte z listy w porównaniu do poprzednich próśb.} \\
\\ 
\texttt{szczur, opalony, czołg, siatkówka, wzgórze, miedziogłów, czerwony, płaszcz, zapasy, niebieski, szydło} \\
            }
        }
    }
    & \\ \\

    \theutterance \stepcounter{utterance}  
    & & & \multicolumn{4}{p{0.6\linewidth}}{
        \cellcolor[rgb]{0.9,0.9,0.9}{
            \makecell[{{p{\linewidth}}}]{
                \texttt{\tiny{[GM$\langle$P2]}}
                \texttt{ZGADNIĘCIE: szydło} \\
            }
        }
    }
     \\ \\

    \theutterance \stepcounter{utterance}  
    & & & \multicolumn{2}{p{0.3\linewidth}}{
        \cellcolor[rgb]{0.9,0.9,0.9}{
            \makecell[{{p{\linewidth}}}]{
                \texttt{\tiny{[GM$|$GM]}}
                ['szydło']
            }
        }
    }
    & & \\ \\

    \theutterance \stepcounter{utterance}  
    & & & \multicolumn{2}{p{0.3\linewidth}}{
        \cellcolor[rgb]{0.9,0.9,0.9}{
            \makecell[{{p{\linewidth}}}]{
                \texttt{\tiny{[GM$|$GM]}}
                {'word': 'szydło', 'assignment': 'team'}
            }
        }
    }
    & & \\ \\

    \theutterance \stepcounter{utterance}  
    & & & \multicolumn{2}{p{0.3\linewidth}}{
        \cellcolor[rgb]{0.9,0.9,0.9}{
            \makecell[{{p{\linewidth}}}]{
                \texttt{\tiny{[GM$|$GM]}}
                {'word': 'szydło', 'assignment': 'team'}
            }
        }
    }
    & & \\ \\

    \theutterance \stepcounter{utterance}  
    & & & \multicolumn{2}{p{0.3\linewidth}}{
        \cellcolor[rgb]{0.9,0.9,0.9}{
            \makecell[{{p{\linewidth}}}]{
                \texttt{\tiny{[GM$|$GM]}}
                {'word': 'płaszcz', 'assignment': 'opponent'}
            }
        }
    }
    & & \\ \\

    \theutterance \stepcounter{utterance}  
    & & & \multicolumn{2}{p{0.3\linewidth}}{
        \cellcolor[rgb]{0.9,0.9,0.9}{
            \makecell[{{p{\linewidth}}}]{
                \texttt{\tiny{[GM$|$GM]}}
                {'hidden': {'team': ['czerwony', 'niebieski', 'opalony'], 'innocent': ['miedziogłów', 'siatkówka', 'czołg', 'szczur', 'zapasy'], 'opponent': [], 'assassin': ['wzgórze']}, 'revealed': {'team': {'team': ['wiertło', 'imam', 'kaznodzieja', 'ksiądz', 'młotek', 'szydło'], 'innocent': ['czarny', 'kij'], 'opponent': ['szczypce', 't-shirt'], 'assassin': []}, 'opponent': {'team': [], 'innocent': [], 'opponent': ['trzepaczka', 'sweter', 'koszula', 'durszlak', 'talerz', 'płaszcz'], 'assassin': []}}}
            }
        }
    }
    & & \\ \\

    \theutterance \stepcounter{utterance}  
    & & \multicolumn{4}{p{0.6\linewidth}}{
        \cellcolor[rgb]{0.9,0.9,0.9}{
            \makecell[{{p{\linewidth}}}]{
                \texttt{\tiny{[P1$\langle$GM]}}
                \texttt{Słowo szydło zostało odgadnięte prawidłowo. Tura twojego kolegi z drużyny się tam skończyła.} \\
\texttt{Teraz podaj kolejną wskazówkę odnoszącą się do niektórych z pozostałych słów twojego zespołu oraz listę powiązanych słów docelowych. Pamiętaj, aby rozpocząć wskazówkę od 'WSKAZÓWKA: ', wstawić nową linię i rozpocząć listę słów docelowych oddzielonych przecinkami od 'CELE: '. Uwaga: niektóre słowa zostały usunięte z list w porównaniu do poprzednich żądań.} \\
\\ 
\texttt{Twoje pozostałe słowa zespołu to: czerwony, niebieski, opalony.} \\
\texttt{Pozostałe słowa twojego przeciwnika to: płaszcz.} \\
\texttt{Pozostałe słowa rozpraszające to: miedziogłów, siatkówka, czołg, szczur, zapasy.} \\
\texttt{Pozostałe słowa zabójcy to: wzgórze.} \\
            }
        }
    }
    & \\ \\

    \theutterance \stepcounter{utterance}  
    & \multicolumn{4}{p{0.6\linewidth}}{
        \cellcolor[rgb]{0.9,0.9,0.9}{
            \makecell[{{p{\linewidth}}}]{
                \texttt{\tiny{[P1$\rangle$GM]}}
                \texttt{WSKAZÓWKA: kolor} \\
\texttt{CELE: czerwony, niebieski} \\
            }
        }
    }
    & & \\ \\

    \theutterance \stepcounter{utterance}  
    & & & \multicolumn{2}{p{0.3\linewidth}}{
        \cellcolor[rgb]{0.9,0.9,0.9}{
            \makecell[{{p{\linewidth}}}]{
                \texttt{\tiny{[GM$|$GM]}}
                \texttt{kolor} \\
            }
        }
    }
    & & \\ \\

    \theutterance \stepcounter{utterance}  
    & & & \multicolumn{2}{p{0.3\linewidth}}{
        \cellcolor[rgb]{0.9,0.9,0.9}{
            \makecell[{{p{\linewidth}}}]{
                \texttt{\tiny{[GM$|$GM]}}
                ['czerwony', 'niebieski']
            }
        }
    }
    & & \\ \\

    \theutterance \stepcounter{utterance}  
    & & & \multicolumn{2}{p{0.3\linewidth}}{
        \cellcolor[rgb]{0.9,0.9,0.9}{
            \makecell[{{p{\linewidth}}}]{
                \texttt{\tiny{[GM$|$GM]}}
                {'word': 'czerwony', 'assignment': 'team'}
            }
        }
    }
    & & \\ \\

    \theutterance \stepcounter{utterance}  
    & & & \multicolumn{2}{p{0.3\linewidth}}{
        \cellcolor[rgb]{0.9,0.9,0.9}{
            \makecell[{{p{\linewidth}}}]{
                \texttt{\tiny{[GM$|$GM]}}
                {'word': 'niebieski', 'assignment': 'team'}
            }
        }
    }
    & & \\ \\

    \theutterance \stepcounter{utterance}  
    & & & \multicolumn{2}{p{0.3\linewidth}}{
        \cellcolor[rgb]{0.9,0.9,0.9}{
            \makecell[{{p{\linewidth}}}]{
                \texttt{\tiny{[GM$|$GM]}}
                \texttt{opponent has won} \\
            }
        }
    }
    & & \\ \\

\end{supertabular}
}

\end{document}
